% FAC-EN-U0059
% Source: fac_body.tex lines 2133--2193.
% Authority: PDF physical pp. 58--59; printed pp. 253--254
% (Chapter III section 3 opening and no. 61 complete).
% French controls: T0034 repairs three inherited epsilon commands to the
% printed membership sign. C0036 preserves printed i_i diplomatically and
% restores the type-consistent terminal index i_q in corrected French/English.

\phantomsection\label{fac:III:3}
\section*{\S3. Cohomology of projective space with values in a coherent
algebraic sheaf}
\addcontentsline{toc}{section}{Chapter III, \S3. Cohomology of projective
space with values in a coherent algebraic sheaf}

\phantomsection\label{fac:III:3:no:61}
\textbf{61. The complexes $C_k(M)$ and $C(M)$.}
We retain the notation of \hyperref[fac:III:1:no:51]{nos.~51} and
\hyperref[fac:III:2:no:56]{56}. In particular, $I$ will denote the
interval $\{0,1,\ldots,r\}$, and $S$ will denote the graded algebra
$K[t_0,\ldots,t_r]$.

Let $M$ be a graded $S$-module, and let $k$ and $q$ be two integers
$\geq0$. We shall define a group $C^q_k(M)$. An element of $C^q_k(M)$ is
a mapping
\phantomsection\label{fac:III:3:no:61:formula:cochain-value}
\[
(i_0,\ldots,i_q)\longmapsto m\langle i_0\cdots i_q\rangle
\]
that assigns to every sequence $(i_0,\ldots,i_q)$ of $q+1$ elements of
$I$ a homogeneous element of $M$ of degree $k(q+1)$ and is alternating
in $i_0,\ldots,i_q$. In particular,
$m\langle i_0\cdots i_q\rangle=0$ if two of the indices
$i_0,\ldots,i_q$ are equal. Addition in $C^q_k(M)$ and multiplication by
an element $\lambda\in K$ are defined in the evident way, and $C^q_k(M)$
is a \textit{vector space over $K$}.

If $m$ is an element of $C^q_k(M)$, define $dm\in C^{q+1}_k(M)$ by
\phantomsection\label{fac:III:3:no:61:formula:differential}
\[
(dm)\langle i_0\cdots i_{q+1}\rangle
=\sum_{j=0}^{j=q+1}(-1)^j t_{i_j}^k.
  m\langle i_0\cdots\widehat{\imath}_j\cdots i_{q+1}\rangle.
\]
\phantomsection\label{fac:III:3:no:61:formula:complex-Ck}
A direct calculation shows that $d\circ d=0$. Hence the direct sum
$C_k(M)=\sum_{q=0}^{q=r}C^q_k(M)$, equipped with the coboundary operator
$d$, is a \textit{complex}; its $q$-th cohomology group will be denoted
by $H^q_k(M)$.

(Following [11], let us mention another interpretation of the elements
of $C^q_k(M)$. Introduce $r+1$ differential symbols
$dx_0,\ldots,dx_r$, and associate with every $m\in C^q_k(M)$ the
``differential form'' of degree $q+1$
\phantomsection\label{fac:III:3:no:61:formula:omega-m}
\[
\omega_m=\sum_{i_0<\cdots<i_q}
m\langle i_0\cdots i_q\rangle\,
dx_{i_0}\wedge\cdots\wedge dx_{i_q}.
\]
\phantomsection\label{fac:III:3:no:61:formula:alpha-k}
If $\alpha_k=\sum_{i=0}^{i=r}t_i^k\,dx_i$, then
\phantomsection\label{fac:III:3:no:61:formula:omega-dm}
\[
\omega_{dm}=\alpha_k\wedge\omega_m;
\]
in other words, the coboundary operation becomes exterior
multiplication by the form $\alpha_k$.)

If $h$ is an integer $\geq k$, let
$\rho^h_k:C^q_k(M)\to C^q_h(M)$ be the homomorphism defined by
\phantomsection\label{fac:III:3:no:61:formula:rho}
\[
\rho^h_k(m)\langle i_0\cdots i_q\rangle
=(t_{i_0}\cdots t_{i_q})^{h-k}
m\langle i_0\cdots i_q\rangle.
\]
We have $\rho^h_k\circ d=d\circ\rho^h_k$, and
$\rho^l_h\circ\rho^h_k=\rho^l_k$ if $k\leq h\leq l$. We may therefore
\phantomsection\label{fac:III:3:no:61:formula:direct-limit-complex}
define the complex $C(M)$, the direct limit of the system
$(C_k(M),\rho^h_k)$ as $k\to+\infty$.
The cohomology groups of this complex will be denoted by $H^q(M)$.
Since cohomology commutes with direct limits (cf. [6], Chapter V,
Proposition~9.3*),
\phantomsection\label{fac:III:3:no:61:formula:direct-limit-cohomology}
\[
H^q(M)=\varinjlim_{k\to\infty}H^q_k(M).
\]

Every homomorphism $\varphi:M\to M'$ defines a homomorphism
\phantomsection\label{fac:III:3:no:61:formula:functorial-chain-map}
\[
\varphi:C_k(M)\longrightarrow C_k(M')
\]
by the formula
$\varphi(m)\langle i_0\cdots i_q\rangle
=\varphi(m\langle i_0\cdots i_q\rangle)$, and hence, by passage to the
limit, a homomorphism $\varphi:C(M)\to C(M')$. Moreover, these
homomorphisms commute with the coboundary and therefore define
homomorphisms
\phantomsection\label{fac:III:3:no:61:formula:functorial-cohomology}
\[
\varphi:H^q_k(M)\longrightarrow H^q_k(M')
\qquad\text{and}\qquad
\varphi:H^q(M)\longrightarrow H^q(M').
\]
If
$0\to M\to M'\to M''\to0$ is exact, then so is the sequence of
complexes
$0\to C_k(M)\to C_k(M')\to C_k(M'')\to0$, and hence there is an exact
cohomology sequence
\phantomsection\label{fac:III:3:no:61:formula:long-exact-sequence}
\[
\cdots\longrightarrow H^q_k(M')\longrightarrow H^q_k(M'')
\longrightarrow H^{q+1}_k(M)\longrightarrow H^{q+1}_k(M')
\longrightarrow\cdots.
\]
The same results hold for $C(M)$ and the $H^q(M)$.

\phantomsection\label{fac:III:3:no:61:remark:no:69}
\textsc{Remark.} We shall see later
(cf. \hyperref[fac:III:4:no:69]{no.~69}) that the $H^q_k(M)$
can be expressed in terms of the $\operatorname{Ext}^q_S$.
